% 中文论文正文；使用 XeLaTeX 编译。
% English title: ContextWorld: Benchmarking In-Context Dynamics Adaptation in Latent Visual World Models
% Uses the official CVPR 2026 style; not a submission-format certificate.
\documentclass[10pt,twocolumn,letterpaper]{article}
\usepackage[pagenumbers]{cvpr}
\usepackage[UTF8,scheme=plain,fontset=fandol]{ctex}
% Shared packages; the official cvpr.sty is kept unchanged.
\usepackage{tabularx}
\definecolor{cvprblue}{rgb}{0.21,0.49,0.74}
\usepackage[pagebackref,breaklinks,colorlinks,allcolors=cvprblue]{hyperref}
\renewcommand{\refname}{参考文献}
\renewcommand{\tablename}{表}
\renewcommand{\figurename}{图}
\patchcmd{\abstract}{Abstract}{摘要}{}{\errmessage{Abstract heading patch failed}}
\def\paperID{*****}
\def\confName{CVPR}
\def\confYear{2027}
\title{ContextWorld：潜在空间视觉世界模型的上下文动力学适应基准}
\author{}
\begin{document}
\maketitle
\begin{abstract}
视觉世界模型学习环境状态随动作变化的规律，即环境动力学，并根据历史视觉观测和后续动作预测未来状态。其中，潜在空间模型将视觉观测编码为潜在状态，并预测其变化。动力学变化时，相同观测与动作也可能产生不同结果。在相关参数或规则未被直接提供时，历史交互可为判断当前动力学提供线索，但总体预测误差难以说明模型是否正确利用了这些线索。为此，我们提出上下文动力学适应基准 ContextWorld。评测时，我们保持当前观测与后续动作不变，替换交互历史，比较预测的变化及其误差，检验模型是否利用历史，以及能否据此准确预测未来状态。基准涵盖四个视觉交互环境中的九项隐藏动力学任务，并提供独立的训练与评测数据。评测数据来自连续交互轨迹，其中的配对样本提供不同动力学条件下的历史与真实未来。模型仅接收历史视觉观测、其间已执行的动作和给定的后续动作，评测期间不更新权重。我们评测了常用的潜在空间世界模型，发现仅在固定动力学条件下训练时，模型对历史线索的利用总体有限。使用基准训练数据的实验表明，扩充数据在多项任务上改善了历史利用，但收益因任务而异；不同模型各有优势任务，有效训练方案也随模型和任务变化。结果表明，ContextWorld 能够揭示不同模型与训练方案在历史利用和预测准确性上的差异。进一步分析发现，在部分训练方案下，训练损失虽下降，模型却仍未能在训练集上依据历史准确预测未来，提示现有训练目标对上下文动力学适应的约束不足。
\end{abstract}
\section{引言}
\label{sec:introduction}

视觉世界模型为规划与控制提供对动作后果的预测 \cite{ha2018worldmodels}。其中，潜在空间世界模型将视觉观测编码为状态表示，并在该空间中预测动作引起的状态变化 \cite{hafner2019planet,hafner2020dreamer}。然而，即使模型能够准确预测训练环境中的动作后果，也未必能应对动力学条件的变化 \cite{wang2026context}。例如，物体外观相同，受到同样推动后的位移却可能因摩擦不同而变化。此时，仅知道物体当前在哪里并不足够，模型还需要判断当前环境如何响应动作。如果训练数据中的动力学条件始终不变，模型就可能沿用学到的规律进行预测，而不从历史中判断当前的动力学条件。

当实际环境的动力学条件与训练时不同或发生变化时，模型需要根据当前条件调整预测，而不能只沿用训练时学到的固定动力学。历史交互中的动作及其结果，可以为判断当前条件提供依据。已有研究为考察这一能力提供了基础。CARL 将环境条件作为可控变量，用于研究条件变化下的泛化与适应 \cite{benjamins2021carl}；世界模型的上下文学习研究区分了从历史辨认已有环境与学习环境规律，并分析上下文长度和环境多样性如何影响两者 \cite{wang2026context}。ICWM 进一步展示了通过任务执行前的交互，在不更新参数的情况下适应相机视角等系统配置变化的可能性 \cite{wang2026icwm}。本文关注其中的视觉动力学预测问题：模型能否在权重固定时，从历史中判断当前动力学，并据此预测给定动作的后果，即上下文动力学适应。

评价这种能力需要明确历史在预测中起了什么作用。模型可能利用历史估计当前运动状态，却没有据此区分不同动力学；也可能准确预测不同条件共有的状态信息，使整体误差较低。因此，一项评测需要让当前观测和查询动作本身不足以确定未来，并比较不同历史是否使预测朝对应的真实未来变化。为此，我们在配对样本中固定当前观测与后续动作，提供不同动力学条件下的连续交互历史及真实未来。对于同一个真实未来，再比较使用当前条件下的历史与其他条件下的历史时的预测误差。这样的对照同时检验历史是否改变预测，以及这种改变是否有助于准确预测。评分使用各模型自身的潜在状态空间，其误差不直接代表统一的物理预测精度。

据此，我们构建 ContextWorld，提供四个视觉交互环境中的九项任务，以及独立的训练与评测数据。任务设计覆盖动作响应幅度、动作延迟、接触后果和状态转移规则，使评测考察不同类型的动力学信息，而不局限于一种运动规律。历史与未来来自真实连续的模拟器交互；模型输入仅含视觉观测、已执行的动作和给定的后续动作，隐藏条件不直接提供，评测期间不更新权重。多数任务检验新场景下对已见条件的辨认，速度任务另测未见参数值的插值与外推。提供训练数据，则使我们能够进一步研究模型在接触不同动力学后，是否学会通过历史调整预测。

我们选取 LeWM~\cite{maes2026leworldmodel}、PLDM~\cite{sobal2025pldm} 和 DINO-WM~\cite{zhou2025dinowm} 建立基线。仅用固定动力学数据（下称原环境数据）训练的模型，在多数任务上对历史中动力学信息的利用有限。为了考察训练经历对这一表现的影响，我们将基准训练数据与原环境数据混合，比较直接从头训练和原环境预训练后继续训练，并考察数据扩量；对 LeWM 和 PLDM，还比较继续更新与冻结图像编码器的方案。各任务分别训练模型。实验发现，扩量在多项任务上有收益，但有效的训练方式因模型与任务而异：直接混合训练时，LeWM 和 PLDM 在动作延迟任务上得分较高，DINO-WM 在推手移动幅度任务上较高；从相同初始化出发，冻结 PLDM 图像编码器并继续训练其余可训练模块，有利于摩擦任务，却不利于阻尼任务。这些差异说明，训练方式的效果需要在不同动力学任务上检验。

训练后表现不足，还可能有两种不同来源：模型在训练中就未学会所需能力，或已学到的能力未能泛化到新场景。为了区分两者，我们将训练损失与固定训练集样本上的能力表现对应起来，并与独立评测数据比较。例如，在 PLDM 摩擦任务的联合训练中，第五至第十轮记录的预测损失均值下降约 18\%，但模型在训练样本上仍未准确预测两种摩擦条件造成的未来差异。这表明，该配置的不足不能仅由训练与测试场景不同解释；优化训练损失与获得所需动力学能力之间仍有差距。结合训练方案对照，我们将这一发现解释为已考察训练目标与优化过程对历史利用的约束不足，而不将其归为所有任务的共同根因。

本文的贡献包括：

\begin{enumerate}
\item \textbf{ContextWorld 基准。} 构建覆盖四个视觉交互环境、九项隐藏动力学任务的训练与评测数据，提供评测协议和评分方法，检验潜在空间世界模型在不更新权重时对交互历史的利用及预测准确性。
\item \textbf{跨模型、数据规模与训练方案的系统比较。} 比较三个模型家族，展示数据扩量、预训练初始化和编码器更新方式对适应能力的影响，以及这些效果如何随模型与任务变化。
\item \textbf{训练过程中的能力不足分析。} 结合训练集与独立评测结果，发现部分训练方案下，模型在训练损失下降时仍未充分学会利用历史中的动力学信息，说明这些不足在训练阶段已经存在，不能仅归因于测试场景的变化。
\end{enumerate}
\section{ContextWorld 基准}
\label{sec:benchmark}

ContextWorld 要分辨模型是从历史判断了当前动力学，还是主要预测了不同条件共有的画面内容。任务构造使隐藏条件对应的未来可以比较；数据划分检验新场景表现；评分分别衡量未来辨别、完整误差与匹配历史的收益。

\subsection{问题定义}
\label{sec:problem}

评测时模型权重固定，只能使用历史图像、已经执行的动作和接下来给定的动作；隐藏参数与模拟器状态不进入模型输入。记模型参数为 $\theta$，$H_t$ 为包含当前图像的观测历史及相邻图像间已执行的动作，$U_{t:t+\tau-1}$ 为接下来给定的动作序列。模型回答这一查询，预测未来图像的表示：

\[
\hat z_{t+\tau}=f_\theta(H_t,U_{t:t+\tau-1}),\qquad
z_{t+\tau}=e_\theta(o_{t+\tau}).
\]

其中，$o_{t+\tau}$ 是模拟器生成的真实未来图像，$e_\theta$ 是本次评测中固定的模型目标编码器。这里的“表示”指模型编码后的图像特征，可以是向量或特征图。真实未来与预测在同一模型的目标表示中比较；不同模型的表示空间不同，误差不构成统一的跨模型物理精度。隐藏参数、模拟器状态和任务标签只用于数据构造与评分。

历史长度 $h$ 指包含当前帧的观测帧数。相邻观测通常间隔五个物理步；因此，$h=3$ 包含两个过去帧、一个当前帧，以及这三帧之间实际执行的动作。这些已观察到的动作后果提供判断当前动力学的线索。

\subsection{任务设计}
\label{sec:tasks}

我们以 TwoRoom 导航环境 \cite{zhou2025dinowm}、PushT 推动环境 \cite{chi2023diffusionpolicy}、DeepMind Control Suite 的 Reacher 环境 \cite{tassa2018control} 和 OGBench 的 Cube 环境 \cite{park2025ogbench} 为基础，构建九项隐藏动力学任务。这些任务分别改变动作响应的幅度或时序、物体在接触中或脱离接触后的运动，以及物体或场景的转移规则。历史线索可能是位移、动作生效时间或门的通行结果。表~\ref{tab:tasks} 列出每项任务需要判断的后果。

\begin{table*}[t]
\centering
\caption{九项任务覆盖不同的动作后果，用于检验模型能否从历史判断隐藏动力学。历史帧数 $h$ 包含当前帧。}
\label{tab:tasks}
\begin{tabularx}{\textwidth}{@{}llcX@{}}
\toprule
任务 & 环境 & 历史帧数 $h$ & 需要从历史判断的动作后果 \\
\midrule
速度 & TwoRoom & 3 & 相同动作使主体移动多远 \\
推手移动幅度 & PushT & 3 & 相同动作使推手移动较短还是较远 \\
机械臂质量 & Reacher & 3 & 相同力矩如何改变机械臂运动 \\
动作延迟 & TwoRoom & 7 & 动作等待多少步才开始生效 \\
接触摩擦 & PushT & 3 & 接触中的摩擦如何改变物体运动 \\
运动阻尼 & PushT & 3 & 脱离接触后，物体运动衰减得快还是慢 \\
夹爪携带规则 & Cube & 3 & 闭合夹爪能否带动物块 \\
门通行规则 & TwoRoom & 3 & 外观相同的门是否允许通过 \\
传送门出口 & TwoRoom & 3 & 进入同一入口后从哪里离开 \\
\bottomrule
\end{tabularx}
\end{table*}

\textbf{成对查询。} 推手移动幅度、机械臂质量、接触摩擦、运动阻尼、夹爪携带和传送门出口使用两种隐藏动力学设置（下称“条件”）的配对样本。每对固定当前画面和接下来给定的动作（查询动作），分别提供各条件下的真实历史与未来。当前画面和动作无法指出应预测哪个未来；对完整、平衡且目标可区分的配对，忽略历史的确定性模型会在两侧给出相同预测，无法同时选对两种未来。

每条样本从历史到查询后的未来都由模拟器连续生成，中途不重置或改写状态。相同画面仍可能隐藏速度或接触状态差异，因此按任务检查查询时物理状态的匹配范围；动作延迟任务保留真实的待执行动作队列。为减少起点泄露条件的可能，对无法自然共享同一起点的任务，生成时交换两种条件的起点，并检查两种条件下起点图像集合一致。具体约束见\href{https://github.com/Anguo-star/ContextWorld/blob/main/docs/Data_Generation.md}{数据构建说明}。

\textbf{其他查询。} 速度任务在相同当前状态和动作下比较多个参数值对应的历史；动作延迟任务比较不同延迟产生的一步未来；门通行任务在同一历史下辨别可通过与被阻挡的未来。动作延迟仍覆盖 0–10 个物理步；仅在五个物理步的一步预测评分中，延迟 5–10 因产生相同的静止响应而合为一组。因此主分使用 0、1、2、3、4、5–10 六个可区分的响应组。

\textbf{泛化范围。} 多数任务检验模型能否在新场景中辨认训练已包含的隐藏条件。速度任务另测未见参数值的插值与外推；门通行任务的测试集使用训练中未见的门位。这些范围分别报告。

\subsection{数据构建与划分}
\label{sec:data}

各任务保存连续交互中的历史图像、已执行动作、查询动作和真实未来。训练集（Training）用于训练，开发集（Development）用于方法选择与预测分析，测试集（Test）用于方法固定后的结果报告。数据划分按任务检查来源轨迹、场景模板和查询图像是否重复，以支持新场景评测；隐藏参数和配对标识不进入模型输入。

推手移动幅度、接触摩擦、运动阻尼、机械臂质量和传送门出口各包含 32,768 个训练配对；夹爪携带任务包含来自 10,000 条不同来源轨迹的 10,000 个训练配对。六项任务的开发集与测试集各含 256 对。配对数不同于独立来源轨迹数：同一来源可产生多个动作窗口，运动阻尼的前向与镜像样本也可能同源。后续统计抽样将同源样本作为一组。

另外三项任务以参数或场景覆盖描述训练集。速度任务覆盖 32 档训练速度，开发集和测试集对每个参考速度各包含 300 个查询。动作延迟的开发集和测试集各包含 300 个查询场景，每个场景覆盖全部 11 种延迟。门通行任务覆盖 96 个训练门位，开发集覆盖 16 个门位、合计 300 个查询，测试集覆盖 42 个门位、合计 300 个查询。速度档数和门位数表示条件覆盖，并非训练样本数。
\subsection{评测协议与评分方法}
\label{sec:evaluation}

任务主分使用 Test 的一步查询，回答模型是否区分了任务规定的未来；Development 上的扩展动作与时域另用于检查完整预测误差和历史收益。九项任务的判断问题不同，主分分别报告，不计算跨任务总分。模型与目标编码器在评测期间保持固定，输入遵循第~\ref{sec:problem}~节的约束。

\textbf{任务成绩。} 对六项二条件配对任务，主分检查匹配历史下的预测是否更接近对应条件的真实未来。记第 $i$ 对中条件 A、B 的预测为 $\hat z_i^A,\hat z_i^B$，真实未来的目标表示为 $z_i^A,z_i^B$；$d$ 为该模型目标表示中的均方误差。对 $N$ 对查询分别判断两侧：

\begin{equation}
\begin{aligned}
A_{\mathrm{future}}&=\frac{100}{2N}\sum_{i=1}^{N}(I_i^A+I_i^B),\\
I_i^A&=\mathbf 1\{d(\hat z_i^A,z_i^A)<d(\hat z_i^A,z_i^B)\},\\
I_i^B&=\mathbf 1\{d(\hat z_i^B,z_i^B)<d(\hat z_i^B,z_i^A)\}.
\end{aligned}
\label{eq:future-score}
\end{equation}

速度任务要求匹配历史的预测误差严格小于同组每种竞争历史的误差；动作延迟按一步未来的六个物理响应组等权平均；门通行在同一历史下辨别可通过与被阻挡的未来。速度总表报告未见速度的插值成绩，其余分布另列。二条件任务和速度任务将平局计为失败；同一预测在配对两侧的比较也相关联，故随机猜测的正确率不能视为无历史模型的预期成绩。主分衡量任务规定的未来辨别，不说明响应幅值或较长轨迹是否准确。

\textbf{完整预测误差。} 选对未来的预测仍可能在位置或运动幅度上失准。我们让同一场景的各隐藏条件共用后续动作序列，模型从真实历史出发自由推演，在第 5、10、15、20、25 个物理步比较预测与真实未来；推演途中不补入真实观测。

为解释误差大小，我们用条件均值构造仅供评分的参照：在给定场景、动作和时刻，对所有条件输出同一未来时，条件等权下的最小平方误差由各条件真实未来的均值取得。对场景 $q$，记该均值为 $\bar z_{qct}=K_q^{-1}\sum_k z_{qkct}$，其中 $K_q,C_q$ 分别是条件数和共用动作候选数，$T=5$ 是评分时刻数；$z_{qkct}$ 和 $\hat z_{qkct}$ 是条件 $k$、候选 $c$、时刻 $t$ 的真实与预测表示。先在场景内等权平均，再汇总场景：

\begin{equation}
\begin{aligned}
L_q&=\frac{1}{K_qC_qT}\sum_{k,c,t}\|\hat z_{qkct}-z_{qkct}\|^2,\\
B_q&=\frac{1}{K_qC_qT}\sum_{k,c,t}\|z_{qkct}-\bar z_{qct}\|^2,\\
E/B&=\frac{\sum_q L_q}{\sum_q B_q}.
\end{aligned}
\label{eq:complete-error}
\end{equation}

$E/B$ 越低越好：0 表示在该目标表示中预测准确，1 表示误差等于条件均值参照，大于 1 表示误差更大。参照只由评分器根据真实未来计算，分母须为正。对条件等权且当前画面、动作相同的查询，始终给出同一未来的确定性无历史模型无法取得小于 1 的比值。这个比值衡量完整误差，不判断匹配历史是否比错配历史有帮助。

\textbf{历史对照。} 保持当前画面、后续动作和真实目标相同，只把匹配当前条件的历史换成同一场景其他条件的历史。History（\%）报告匹配历史使预测误差严格更小的比例；速度任务也将此胜率用作主分。

胜率不显示误差改善的幅度。令 $L_q^+=L_q$ 为匹配历史下的场景误差，$L_q^-$ 为错配历史下的误差；有多个错配条件时对其等权平均。沿用相同的条件、候选和时刻权重，记 $E_+=\sum_qL_q^+$、$E_-=\sum_qL_q^-$、$B=\sum_qB_q>0$，计算

\begin{equation}
G_H=\frac{E_--E_+}{B}.
\label{eq:history-gain}
\end{equation}

$G_H>0$ 表示匹配历史降低了误差，改善后的完整误差仍由 $E/B$ 描述。

\textbf{条件响应诊断。} 总体预测损失可能因为两种条件共有的偏差减小而下降。训练分析另比较两条件造成的未来差异，即条件响应。对二条件查询，令真实差异为 $\Delta z_i=z_i^B-z_i^A$，预测差异为 $\Delta\hat z_i=\hat z_i^B-\hat z_i^A$；归一化响应误差（NRE）衡量两者是否一致：

\begin{equation}
\mathrm{NRE}=\frac{\sum_i\|\Delta\hat z_i-\Delta z_i\|^2}
{\sum_i\|\Delta z_i\|^2}.
\label{eq:response-error}
\end{equation}

NRE 越低越好；0 表示条件响应准确，1 是预测完全不随条件变化时的误差水平。取条件差会消去共有的预测偏差，所以 NRE 不能代替完整误差，单独的 NRE=1 也不能证明模型忽略历史。训练集与 Development 的诊断分别报告预测时刻；NRE 仅在分母为正时定义。

\textbf{结果比较。} 任务主分对应各任务规定的辨别问题。E/B 与 NRE 在各检查点自身的目标表示中计算，不能直接用于跨模型的物理精度排名。若目标编码器随训练改变，跨检查点的表示误差比较也受到表示变化影响；训练诊断只描述各检查点在自身目标表示下的表现。同一目标表示中的误差对照更直接。这些读数描述同一组预测的不同方面，不作为相互独立的多份证据。

相同配置的训练结果等权汇总，标准差反映训练间波动。多步误差的区间按来源轨迹或场景组进行 bootstrap 重采样 \cite{efron1993bootstrap}，反映给定检查点下的场景抽样不确定性。原环境中的交叉熵法（CEM）\cite{deboer2005cem} 规划成绩另列，用于检查原有规划能力是否保留；原环境任务未必要求辨认变化后的动力学，故该成绩不替代适应能力评测。
{\small
\bibliographystyle{ieeenat_fullname}
\bibliography{references}
}
\end{document}

% 以下图注在对应图表制作后插入正文。
% Manuscript figure and result-table captions. Insert with their matching assets.

% ### 图 1　从交互历史到隐藏动力学预测
% \textbf{ContextWorld 的受控查询与任务范围。} 接触摩擦示例中，两种隐藏条件的真实历史到达相同当前画面；模型依据历史，在相同后续动作下预测各自未来。模型只见图像和动作，真实未来经该模型固定的目标编码器转换后用于评分。九项任务覆盖四个环境；图中配对结构对应六项二条件任务，其余任务使用各自的查询构造。

% ### 表 2　三个模型家族在九任务 Test 上的上下文适应表现
% \textbf{九任务 Test 主指标（\%，越高越好）。} T0 仅用原环境数据训练；T1 直接以原环境与任务数据按 50/50 混合训练；T2 从 T0 权重出发，以同样混合数据联合更新；T3 从相同 T0 权重出发，仅冻结图像编码器，其余模块继续训练。DINO-WM 的预训练视觉编码器始终固定，故无 T2。表中 86 个模型／方案／任务组合来自各任务分别训练的检查点，各格为相同配置训练重复的均值。速度列报告未见速度插值，延迟列使用七帧历史；“—”表示未设可比检查点。各任务按自身规则评分，不合成总分；完整预测误差与原环境规划成绩另报。

% ### 图 2　冻结图像编码器的效果随任务改变
% \textbf{PLDM 相同初始化下更新或冻结图像编码器的 Test 对照。} 两方案共用原环境权重，并使用相同任务的 32,768 个配对、50/50 混合比例、10 轮第二阶段训练和批量设置；仅图像编码器是否更新不同，预测器等其余可训练模块继续训练。冻结后，接触摩擦任务主指标由 50.00\% 升至 73.63\%，运动阻尼由 79.88\% 降至 54.49\%。两任务反转表明冻结收益依赖任务，不能单独归因于某条梯度路径。

% ### 图 3　完整预测改善可以没有相应的条件响应改善
% \textbf{DINO-WM 的完整误差与条件响应。} 推手移动幅度、接触摩擦和运动阻尼各使用相同的 256 个 Development 场景、候选动作及第 5、10、15、20、25 个物理步输出；T0、T1 共用冻结视觉编码器。E/B 为完整预测误差与条件均值参照误差之比，NRE 衡量预测与真实条件响应之差，1 为零响应参照。阻尼的 E/B 由 9.284 降至 7.594，NRE 由 1.001 变为 1.068；完整误差降低后，条件响应误差仍接近零响应基线的误差水平。完整误差还包含两条件共有的预测偏差，两项读数来自同组输出，不是独立证据。区间按来源组 bootstrap 重采样，表示给定检查点的来源抽样不确定性；误差仅在模型自身表示中比较。历史收益另由匹配与错配历史对照检验。

% ### 图 4　训练损失下降不保证训练查询上的条件响应改善
% \textbf{PLDM 摩擦与阻尼的训练损失及条件响应。} 第 5–10 轮曲线显示已记录训练批次的均值、中位数及批次四分位区间；NRE 只在第 5、10 轮保存的检查点测量，使用固定 512 对训练查询和全部 256 对 Development 查询，预测当前状态后第 5 个物理步。NRE=0 为准确条件响应，1 为零响应参照。摩擦联合训练的总损失、预测损失均值分别下降 9.37\%、18.05\%，训练 NRE 为 1.009→1.011；冻结图像编码器后为 0.674→0.587。阻尼的联合训练 NRE 为 0.236→0.197，冻结训练为 0.940→0.867；Development 也呈摩擦更有利于冻结、阻尼更有利于联合更新的方向。损失与 NRE 来自不同样本集合，各检查点使用自身目标表示；两个 NRE 检查点不构成完整能力曲线。
